% Related work section for OrbitGraph paper.

\section{Related Work}
\label{sec:related}

Prior LEO routing work spans snapshot and centralized schemes, distributed
protocols, satellite SDN management, evaluation platforms, and broader
networking literature, which we contrast with OrbitGraph's graph-centric,
proactive SDN forwarding model.

\subsection{Centralized and snapshot routing}

LEO constellations have been modeled as finite-state automata whose states are
recurring visibility configurations, with inter-satellite link assignment and
multi-commodity routing optimized offline per
state~\cite{seong:topological-design:1995}.
That strategy avoids transient instability but relies on offline computation for
a small constellation (12 satellites) and ignores online failure recovery.
OrbitGraph shares the snapshot view but computes shortest paths online and
installs kernel routes through an SDN controller rather than storing per-state
tables on satellites.

Wide-area routing in mega-constellations with and without inter-satellite links
(ISLs) has been studied, emphasizing ground-relay paths when ISLs are absent and
hierarchical, tile-based updates when centralized Dijkstra over thousands of
relays is too slow~\cite{handley:ground-relays:2019}.
Unlike OrbitGraph, that line targets constellation design and relay selection,
not in-container forwarding during live GSL handover, and does not compare
against a distributed interior gateway protocol (IGP) baseline.

\subsection{Distributed LEO routing protocols}

DisCoRoute is a distributed on-demand algorithm for Walker-Delta
mega-constellations that derives a closed-form minimum hop count between orbital
identifiers and selects the next hop in constant time, reporting large speedups
at up to 32\,000 satellites~\cite{gregory:satellite-routing:2022}.
It optimizes hop count on a fixed +Grid ISL pattern and does not model link
failures or convergence delay.
OrbitGraph instead centralizes shortest-path computation on delay-weighted
graphs and measures end-to-end outage during handover.

Orbit-predictive shortest path first (OPSPF) keeps OSPF-style hello and
link-state machinery but uses orbit prediction so regular motion generates no
protocol traffic, with only failures triggering updates~\cite{pam:opspf:2019}.
Evaluated on 48 satellites, it reduces convergence time over vanilla OSPF.
OrbitGraph removes distributed flooding entirely, recomputing $G_t$ from
ephemeris and programming routes directly to target zero handover outage through
proactive install ordering.

Our OSPF baseline, BIRD on StarryNet~\cite{lai2023starrynet}, represents the
reactive link-state approach from early protocol studies~\cite{sidhu:ospf:1993}
and widely used in LEO evaluations~\cite{manzanareslopez2025comprehensive}.
It reconverges only after full topology mutation, which
Section~\ref{sec:analysis} shows can black-hole traffic for several seconds at
GSL handover, the behavior OrbitGraph eliminates.

Terrestrial ad hoc routing, such as ad hoc on-demand distance vector
(AODV)~\cite{perkins:aodv:1999} and greedy perimeter stateless routing
(GPSR)~\cite{karp:gpsr:2000}, targets unpredictable mobility without orbital
models, and ns-3-leo reports that such mobile ad hoc network (MANET) protocols
perform poorly in dynamic satellite scenarios~\cite{ns3leo2022}.
OrbitGraph exploits the opposite assumption: topology is predictable from
SGP4-based delay matrices, so a centralized graph controller is feasible.

A recent robustness study shows that inter-satellite link failures from
environmental and load factors degrade static routing far more than dynamic or
distributed approaches~\cite{yang:leo-routing-resilience:2025}, motivating the
failure-mode analysis we report for OrbitGraph under satellite damage.

\subsection{SDN and graph-based network management}

The SDN paradigm separates control and data planes and enables programmatic
forwarding through interfaces such as OpenFlow~\cite{10.1145/1355734.1355746}.
Representing an SDN domain as a graph gives controllers a unified structure for
topology-aware management and shortest-path algorithms~\cite{7014202}, which
OrbitGraph applies to time-varying LEO snapshots rather than terrestrial
OpenFlow switches.

Satellite-specific SDN research has focused on controller placement, not
data-plane route computation.
Dynamic placement on an Iridium-class 66-satellite constellation has been cast
as an integer linear program minimizing average flow setup
time~\cite{papa:sdn-cpp:2018}, while static placement adapts
switch-to-controller bindings to topology snapshots on Iridium and Celestri
models~\cite{guo:spda:2022}.
Both optimize control placement and treat forwarding as a black box, whereas
OrbitGraph addresses what routes to install and when relative to link lifetime,
evaluating data-plane availability during handover rather than switch-controller
latency.

Centralized traffic engineering has also been applied to satellite backbones,
where a controller computes low-latency routes across a time-varying topology to
meet flow-level objectives~\cite{wu:sate-satellite-te:2025}.
That work optimizes path selection and load balancing at the traffic-matrix
level, whereas OrbitGraph targets per-handover forwarding continuity measured as
data-plane outage on emulated nodes.

\subsection{Simulation and emulation platforms}

The ns-3-leo module~\cite{ns3leo2022} adds orbital mobility, ISLs, and MANET
routing to discrete-event simulation, highlighting limitations of terrestrial
protocols in space.
A dedicated LEO simulator provides configurable +Grid and sparse ISL topologies
with 3D visualization~\cite{kempton:network-simulator:2021}, and the underlying
SILLEO-SCNS tool~\cite{kempton2020silleo,online:silleo:2025} focuses on topology
generation rather than packet-level stacks.
Simplified OMNeT++~5 scenarios support rapid
prototyping~\cite{dewangan:simple-omnet:2017}.
These platforms either lack real protocol stacks or use distributed routing
without centralized proactive install.

Packet-level ns-3 frameworks are surveyed
comprehensively~\cite{manzanareslopez2025comprehensive}.
Hypatia~\cite{kassing2020exploring} runs ns-3 with static Floyd--Warshall
routing, StarPerf~\cite{lai2020starperf} profiles area-to-area performance, and
LEOPath~\cite{leopath2024} targets routing-algorithm comparison without
end-to-end kernel forwarding.
StarryNet~\cite{lai2023starrynet} emulates real TCP/IP and OSPF in containers
and is our evaluation harness, which OrbitGraph extends with a centralized
controller, incremental kernel route install, and proactive handover
(Section~\ref{sec:simulation}).
A recent review of LEO/NTN routing spanning classical, learning-based, and
quantum methods finds the evidence base predominantly simulation-driven, with
limited testbed or emulation support and often simplified control-plane
realism~\cite{mattar:leo-ntn-routing-perspectives:2026}, a gap OrbitGraph
addresses by evaluating real kernel forwarding under a centralized controller in
container emulation.

\subsection{Broader LEO networking literature}

Growing interest in large LEO and swarm systems motivates global, low-latency
connectivity~\cite{smallsats2020,asri2024swarm,farrag:swarm-survey:2021}.
Swarm intelligence frameworks~\cite{chen2022swarm} and formation control
studies~\cite{fangchen2022pdg} address mission-layer coordination, not IP
handover, and machine-learning routing is surveyed with emphasis on model design
and deployment~\cite{zhu:ml-routing-survey:2025}, with recent work applying deep
graph reinforcement learning to routing on large-scale LEO
topologies~\cite{huang:graph-rl-leo-routing:2026}.
An earlier environment for distributed routing analysis in LEO
swarms~\cite{online:silleo-routing:2025} compared conventional protocols without
a centralized proactive strategy.
OrbitGraph instead uses explicit graph snapshots and Dijkstra shortest paths
with deterministic SDN install ordering, quantifying data-plane outage on the
same class of emulated constellations for reproducible comparison with OSPF.

\subsection{Positioning of OrbitGraph}

Table~\ref{tab:related-positioning} summarizes how OrbitGraph relates to the
closest prior lines.
Prior snapshot and centralized schemes optimize offline tables or relay
placement, distributed protocols optimize hop count or OSPF convergence, and
satellite SDN papers optimize controller placement.
OrbitGraph combines online graph snapshots, centralized delay-weighted shortest
paths, incremental dataplane deltas, and proactive GSL handover in one emulated
system with measured end-to-end outage, a combination absent from these works.

\begin{table}[!h]
  \centering
  \caption{OrbitGraph vs.\ representative related work (selected dimensions).}
  \label{tab:related-positioning}
  \renewcommand{\arraystretch}{1.1}
  \footnotesize
  \begin{tabular}{@{}lcccc@{}}
    \hline
    Work & Central & Real FIB & Proactive HO & Outage eval. \\
    \hline
    Snapshot routing~\cite{seong:topological-design:1995} & Yes & No  & offline & No \\
    DisCoRoute~\cite{gregory:satellite-routing:2022} & No  & No  & No      & No \\
    OPSPF~\cite{pam:opspf:2019} & No  & Yes & partial & partial \\
    SDN placement~\cite{papa:sdn-cpp:2018,guo:spda:2022} & Yes & No  & No      & No \\
    StarryNet OSPF~\cite{lai2023starrynet} & No  & Yes & No      & Yes \\
    OrbitGraph (this work) & Yes & Yes & Yes     & Yes \\
    \hline
  \end{tabular}
  \\[0.4em]
  \footnotesize
  HO: handover. Real FIB: kernel or protocol forwarding on emulated nodes.
\end{table}
